% Generated draft: requires mathematical and source review.
\documentclass[11pt]{article}
\usepackage[margin=25mm]{geometry}
\usepackage{fontspec}
\setmainfont{Latin Modern Roman}
\usepackage{amsmath,amssymb,mathtools,unicode-math}
\setmathfont{Latin Modern Math}
\usepackage{enumitem,needspace,xurl,hyperref}
\setlength{\parindent}{0pt}
\setlength{\parskip}{5pt}
\setlength{\emergencystretch}{3em}
\newcommand{\N}{\mathbb{N}}
\newcommand{\Z}{\mathbb{Z}}
\newcommand{\Q}{\mathbb{Q}}
\newcommand{\R}{\mathbb{R}}
\newcommand{\C}{\mathbb{C}}
\newcommand{\F}{\mathbb{F}}
\newcommand{\A}{\mathbb{A}}
\newcommand{\PP}{\mathbb{P}}
\newcommand{\E}{\mathbb{E}}
\newcommand{\eps}{\varepsilon}
\newcommand{\dd}{\,\mathrm{d}}
\begin{document}
\section*{UnsolvedMath ID: 701221}
\textbf{Source problem code:} catalogue entry
\subsection{Definitions and mathematical statement}
Let \(G\) be a finite simple group of Lie type, with a specified defining characteristic \(p\). Interpret a \emph{semisimple abelian subgroup} as an abelian subgroup all of whose elements are semisimple; in this finite setting, this means an abelian subgroup of order relatively prime to \(p\). Set
\[
\mathcal{S}_p(G)=\{A\leq G: A\text{ is abelian and }p\nmid |A|\}.
\]
For \(x\in G\), write \(A^x=x^{-1}Ax\). Here \(H<G\) means that \(H\) is a proper subgroup, and \(1\) denotes the trivial subgroup.

Determine which finite simple groups of Lie type satisfy
\[
\forall A\in\mathcal{S}_p(G)\quad
\forall H<G\quad
\exists x\in G\quad
A^x\cap H=1.
\]
The conjugating element \(x\) may depend on both \(A\) and \(H\).

The supplied source reports that the required conjugate disjointness holds when \(A\) is cyclic, citing J. Siemons and A. Zalesskii, \emph{J. Algebra} \textbf{256} (2002), 611--625. It also reports that it holds when \(G=\mathrm{PSL}_n(q)\) and \(A\) is contained in some maximal torus, citing the same authors, \emph{J. Algebra} \textbf{226} (2000), 451--478. These are supplied background assertions, not independently verified here.

The source gives the following obstruction, communicated to the proposer by V. I. Zenkov: for \(G=L_2(5)=\mathrm{PSL}_2(5)\), there are subgroups \(A\cong C_2\times C_2\) and \(H\cong C_5\rtimes C_2\), denoted \(2\times2\) and \(5:2\) in Atlas notation, such that
\[
\forall x\in G,\qquad A^x\cap H\ne1.
\]
Here \(C_m\) denotes the cyclic group of order \(m\), and \(5:2\) is the dihedral group of order \(10\). With the displayed defining characteristic \(5\), this \(A\) is semisimple abelian.

The source does not explicitly define ``semisimple abelian.'' The elementwise interpretation above is used for this draft; it does not impose containment in a maximal torus. The defining realization must be retained when interpreting the question, since exceptional isomorphisms can identify groups presented in different defining characteristics. The source does not specify how such presentations should be treated in the requested classification.
\subsection{Short English statement}
Classify the finite simple groups of Lie type for which every semisimple abelian subgroup has a conjugate meeting any prescribed proper subgroup trivially. Semisimplicity is interpreted relative to the defining characteristic. The source supplies positive results under additional restrictions and a counterexample in \(L_2(5)\); its registry status is \texttt{open}, without independent verification here.
\subsection{Sources}
\begin{itemize}
\item Supplied local UnsolvedMath registry record: permanent ID \(701221\), exact source code \texttt{catalogue entry}, title ``Kourovka Notebook Problem 16.29.'' The record attributes the problem to E. P. Vdovin, proposed in 2006, and labels it \texttt{open}. Its recorded update date is \texttt{2026-10-16}; this is registry metadata, not an independent literature-verification date.
\item \emph{The Kourovka Notebook (Unsolved Problems in Group Theory)}, Problem 16.29, originally in the 16th issue (2006); the supplied background also identifies the 21st issue edition (2026), \texttt{arXiv:1401.0300v48}. The supplied collection description names E. I. Khukhro and V. D. Mazurov as editors.
\item J. Siemons and A. Zalesskii, \emph{J. Algebra} \textbf{256} (2002), 611--625. Supplied reference for the cyclic-subgroup assertion; no article title or link is supplied.
\item J. Siemons and A. Zalesskii, \emph{J. Algebra} \textbf{226} (2000), 451--478. Supplied reference for the maximal-torus assertion in \(\mathrm{PSL}_n(q)\); no article title or link is supplied.
\item V. I. Zenkov: the supplied statement attributes the \(L_2(5)\) example to a communication to the proposer. No separate publication or communication date is supplied.
\end{itemize}
\end{document}
